Antes que nada quisiera agradecer a la Universidad de Rosario y todos los integrantes del Departamento de Ciencias de la Computación. Cuya participación educativa fue imprescindible para mi desarrollo profesional.

A mi madre y padre, cuyo ejemplo y presencia (y muchas veces contención emocional) fue lo que me permitió perseverar durante tantos años de estudio.

A mí familia extendida, los que ya no están y los que siguen estando. Quienes me permitieron momentáneamente olvidar el estrés con cada comida compartida.

A mis amigos, los cuales se bancaron incontables despotricaciones sobre mi vida académica, y cuya existencia me recordaba que todo esto, al final del día, valía la pena.

A mis compañeros, quienes me siento afortunado de  también incluir en el anterior grupo, pero merecen además una mención especial. No solo compartimos muchas penas y glorias juntos (sumado a una pandemia mundial), sino que además no me es posible exagerar como su apoyo y camaradería fueron los motores que me permitieron llegar hasta acá. A todos ustedes, gracias infinitas.

Al doctor Mario Abadi por su aporte en el capítulo introductorio.

A mí director Juan Cabral, mí co-director Pablo Granitto, y la licenciada Valeria Cristiani, cuyo aporte y guía fueron indispensables para completar este trabajo.

Y por última a Dios, quien me bendijo con todas las herramientas y oportunidades que me brindó, pero aún más importante, rodeó de tantas personas que amo y apoyan en todo lo que me propongo.